% Response to CGTA referee report (25 October 2026).
% Build from the repository root:
%   pdflatex -interaction=nonstopmode response_to_associate_editor.tex
%   pdflatex -interaction=nonstopmode response_to_associate_editor.tex
%   pdflatex -interaction=nonstopmode response_to_associate_editor.tex
\documentclass[11pt]{article}
\usepackage[T1]{fontenc}
\usepackage{graphicx} % Required for inserting images
\usepackage{xcolor}
\usepackage[round]{natbib}
\usepackage{newtxtext}
\usepackage{amsmath,amssymb}
\usepackage[margin=1in]{geometry}

\newcommand{\tr}[1]{\textcolor[HTML]{aa0000}{#1}}
\newcommand{\tb}[1]{\textcolor[HTML]{0000bb}{#1}}
\newcommand{\ta}[1]{\textcolor[HTML]{99007d}{#1}}
\newcommand{\tg}[1]{\textcolor[HTML]{009999}{#1}}

\title{Response to Associate Editor}
\author{Heriberto Espino-Montelongo \and H\'ector Maravillo}
\date{November 3, 2026}

\begin{document}
\maketitle

\noindent Dear Associate Editor,

Regarding our research article entitled ``Geometry and Volume of
Stepping-Stone Diversion Neighborhoods in Euclidean Spaces of Arbitrary
Dimension'' (manuscript number CGTA-D-26-00059), we thank you and the
reviewer for the careful evaluation and constructive suggestions.
%Your comments have been fundamental in significantly improving the overall quality of the manuscript.

We have revised the manuscript accordingly, marking the changes in dark red.
%We have made our best effort to address all of the issues raised in your comments. Below, we provide a detailed, point-by-point response to each of your comments.
Our point-by-point responses follow. The comments are reproduced or
lightly paraphrased where mathematical notation requires clarification.

\begin{enumerate}

\item \tb{Section 4 does a great job of gathering all inclusion and
monotonicity results that computational geometers might be interested in.}

\textbf{Response.} We thank the reviewer for this positive assessment.
We have retained the organization and results of Section~4.

\item \tb{In formula (56), please mention explicitly what the function is
and provide a reference.}

\textbf{Response.} We now identify $I_x(a,b)$ explicitly as the
\emph{regularized incomplete beta function}, give its definition in terms
of the incomplete and complete beta functions, and cite the
hyperspherical-cap formula of Li (2011) in the revised
proposition.

\item \tb{You take $\alpha\geq1$ throughout the paper. Can you also consider
$\alpha<1$? Maybe you could comment on that.}

\textbf{Response.} We added a remark following the special-cases result.
For $0<\alpha<1$, strict subadditivity gives
$r^\alpha+s^\alpha>(r+s)^\alpha$ for $r,s>0$.
Together with the triangle inequality, this shows that the diversion
neighborhood consists only of its two endpoints. Consequently, every
finite-set graph is complete, and its $d$-dimensional normalized volume
is zero. We therefore retain $\alpha\geq1$ as the main domain of study.

\item \tb{The definition in (5) does not directly admit $\alpha=\infty$.
Consider the equivalent expression with the $1/\alpha$ power so that
its limit can be taken explicitly, and discuss the endpoint.}

\textbf{Response.} We have rewritten the finite-parameter defining
inequality in the equivalent form
\[
  \bigl(\lVert z-p\rVert^\alpha+\lVert z-q\rVert^\alpha\bigr)^{1/\alpha}
  \leq \lVert p-q\rVert.
\]
Its left-hand side converges to
$\max\{\lVert z-p\rVert,\lVert z-q\rVert\}$ as $\alpha\to\infty$.
We now distinguish the closed-lune set obtained from the limiting
non-strict inequality from the open-lune relative-neighborhood graph
obtained by eventual stabilization on finite point sets. This
boundary distinction explains why the limiting geometry and the
limiting graph need not coincide in degenerate configurations.

\newpage
\item \tb{Section 5 applies theorems of Devroye and Watanabe to derive
some results on the expected degree and the expected edge length for
random stepping-stone graphs. It has no new technical contributions.
It would have been nice, though, if the results were more homogeneous.
In Devroye's model, the expected edge lengths are missing, and in the
Poisson model for dimension 2, the expected degree results are missing.}

\textbf{Response.} We agree that the two sampling models should be
clearly distinguished. We have added an explicit result for the
stationary planar Poisson model. With
$a=a_{2,\mathrm{SS}}(\alpha)$, its mean degree is $\pi/a$,
consistent with the degree constant of Devroye (1988).
Hence the intensity of undirected edges is $\rho\pi/(2a)$.
Combining this with the mean edge length in
Watanabe (2008) gives the expected total edge length
per unit area,
\[
  \frac{\rho\pi}{2a}\,
  \frac{\sqrt{\pi}}{2\sqrt{\rho a}}
  =\frac{\pi^{3/2}\sqrt{\rho}}{4a^{3/2}}.
\]
The revised manuscript states this result as a proposition with a
short proof. We keep Devroye's general independent-sample model
separate: its cited theorem provides asymptotic degree and edge-count
conclusions, not the corresponding expected-length formula for an
arbitrary density. We have not claimed that such a formula follows
from that theorem.

\end{enumerate}

\noindent We appreciate the reviewer's suggestions and hope that these
clarifications address the concerns raised.

\end{document}
